\documentclass[conference]{IEEEtran}
\usepackage[utf8]{inputenc}
\usepackage{amsmath}
\usepackage{booktabs}
\usepackage{hyperref}
\usepackage{graphicx}

\begin{document}

% --- SECTION A: INFORMATION ---
\section*{A. Information}
\noindent \textbf{Project's Name:} Implicit Feedback-Based Poetry Book Recommendation using Graph-Based Learning \\
\textbf{Project's Code:} A25 – Graph-Based Recommendation \\
\textbf{Course:} Social Network Analysis – IS353.Q21.CTTT \\
\textbf{Duration:} 01/2026 - 05/2026 \\
\textbf{Group:} 13 – Group 36 \\
\textbf{Members:}\\
1. Nguyễn Minh Dũng (23520333) \\
2. Võ Minh Quân (23521274) \\
3. Nguyễn Thị Quỳnh Nhi (23521111) \\
4. Lê Tự Nhân (23521076) \\
\textbf{Supervisor:} Trần Hưng Nghiệp

\vspace{1.5em}

% --- SECTION B: CONTENT ---
\section*{B. Content}

\subsection*{1. Mini Problem Formulation}
Modern Recommender Systems have shifted towards Graph Learning-based Recommender Systems (GLRS). In this project, we model user-book interactions as a bipartite graph where nodes represent users and items, and edges denote implicit feedback. A critical challenge identified in sparse datasets like Goodreads is the "static" treatment of interactions in traditional collaborative filtering. As user preferences are dynamic, older interactions may no longer accurately reflect current interests. We propose integrating temporal signals into the graph propagation mechanism of LightGCN to mitigate data sparsity and capture the shifting behavioral patterns of users.

\subsection*{2. Related Work \& 3. Data}
Our work builds upon foundational Matrix Factorization models and recent advancements in Graph Convolutional Networks (GCNs). Specifically, we leverage the LightGCN architecture, which simplifies the graph convolution process by removing non-linear activations. We utilize the Goodreads poetry subset (provided by the UCSD Recommender Systems Group), which contains detailed interaction logs with timestamps, allowing for the modeling of temporal dynamics. For a detailed overview of the data preprocessing and splitting procedure, please refer to the Data Pipeline in the Appendix (Figure \ref{fig:data_pipeline}).

\subsection*{4. Evaluation Methodology}
The models are evaluated using standard top-$K$ ranking metrics: \textbf{Recall@$K$}, \textbf{NDCG@$K$}, and \textbf{HitRatio@$K$}. 
\begin{itemize}
    \item \textbf{Recall@$K$:} Measures whether the true ground-truth item appears within the top-$K$ recommended list.
    \item \textbf{NDCG@$K$:} Measures ranking quality by considering the specific position of the hit item; higher positions yield higher scores.
    \item \textbf{HitRatio@$K$:} Measures the probability of a successful recommendation (at least one relevant item) within the top-$K$ list.
\end{itemize}
To ensure a robust predictive evaluation, we follow a \textbf{Temporal Leave-One-Out} protocol. In this setup, the most recent interaction for each user is held out for the test set, and the second most recent for validation (see Figure \ref{fig:evaluation_protocol} in the Appendix for the complete evaluation flowchart).

\subsection*{5. Baselines \& Phase 1 Preliminary Results}
We evaluate several foundational and graph-based models, categorized by their implementation status:
\begin{itemize}
    \item \textbf{PopRec}: Recommends the most popular items globally, serving as a measure of popularity bias.
    \item \textbf{Standard MF}: A Latent Factor model optimized using \textbf{BPR Loss}. Both MF and BPR-MF are fully implemented in Phase 1.
    \item \textbf{LightGCN}: A theoretical graph-based baseline currently in the literature review and architectural study phase (no experimental results yet).
\end{itemize}

\subsubsection*{Phase 1 Benchmark Analysis}
Phase 1 focused on implementing and comparing standard latent factor models across two datasets: \textit{ml-100k} (initial benchmark) and our target \textit{Goodreads Poetry} subset. The results for standard MF and BPR-MF are summarized in Table \ref{tab:comparison_results}.

\begin{table}[!t]
    \centering
    \caption{Phase 1 Benchmarks (Average @ K=10, Test Split)}
    \label{tab:comparison_results}
    \begin{tabular}{llccc}
        \toprule
        \textbf{Model} & \textbf{Dataset} & \textbf{Recall@10} & \textbf{NDCG@10} & \textbf{Hit@10} \\
        \midrule
        Standard MF & ml-100k & 0.0855 & 0.0998 & 0.5271 \\
        BPR-MF & Goodreads & 0.7416 & 0.5807 & 0.7416 \\
        \bottomrule
    \end{tabular}
\end{table}

\textbf{Analysis}: The high \textit{HitRatio} in standard MF compared to its \textit{MRR} (0.1915) indicates that while the model identifies correct items, they are often ranked lower in the list due to **Popularity Bias**, where the model favors globally trending items over personalized context. While BPR-MF shows significantly higher metrics on the larger target dataset, it remains a "static" model without temporal awareness.

\textbf{Logic \& Parameters}: The MF models optimize user/item embeddings in a 16-dimensional space using \textbf{BPR Loss}, which employs a pairwise ranking logic to optimize the difference between positive and randomly sampled negative items. The hyperparameters include a Learning Rate of 0.005, a Batch Size of 2048, and a training time of $\approx 0.1$s/epoch.

\subsection*{6. Proposed Method}
We introduce \textbf{Edge-weighted LightGCN with temporal decay}. This model generalizes LightGCN by weighting interaction edges by their recency to address the "stagnant preference" weakness identified in Phase 1:
\begin{equation}
    w_{ui}^{time} = e^{-\lambda \cdot age_{ui}}
\end{equation}
Where $age_{ui} = T_{now} - T_{ui}$. By prioritizing current interests during graph propagation, the model mitigates global popularity bias and captures evolving tastes.

\subsection*{7. Expected Results}
We expect the proposed model to show significant accuracy improvements, particularly in ranking precision (MRR/NDCG), by effectively partitioning popularity noise from semantic user interest.

\subsection*{8. System Prototype Demo}
The system prototype integrates our backend models with a high-performance serving layer to provide real-time recommendations.

\noindent \textbf{Demo Prototype Link:} \url{INSERT_LINK_HERE}

\vspace{2em}

% --- APPENDIX ---
\section*{Appendix}
\subsection*{Experimental Plan}

\begin{table}[!t]
    \centering
    \caption{Dataset Statistics (Goodreads Poetry Subset)}
    \begin{tabular}{lrr}
        \toprule
        \textbf{Metric} & \textbf{Before Filtering} & \textbf{After Filtering} \\
        \midrule
        Users & 377,799 & 191,909 \\
        Books & 36,514 & 36,326 \\
        Interactions & 2,734,350 & 2,482,575 \\
        Sparsity & 0.999802 & 0.999644 \\
        \bottomrule
    \end{tabular}
\end{table}

\begin{table}[!t]
    \centering
    \caption{Ablation Study Variants}
    \begin{tabular}{ll}
        \toprule
        \textbf{Variant} & \textbf{Modification} \\
        \midrule
        Without Temporal Weighting & LightGCN baseline ($w_{ui}^{time} = 1$) \\
        Alternative Temporal Function & Linear vs. Exponential decay comparison \\
        Only Temporal Weighting & No graph propagation layers ($L=0$) \\
        \bottomrule
    \end{tabular}
\end{table}

\begin{figure}[!t]
    \centering
    \includegraphics[width=\linewidth]{data_pipeline.png}
    \caption{Data Pipeline: Raw dataset $\rightarrow$ Filter year $\ge$ 2007 \& Remove users $<$ 3 interactions $\rightarrow$ Group \& Sort interactions by time $\rightarrow$ Map IDs to index $\rightarrow$ Temporal Leave-One-Out split.}
    \label{fig:data_pipeline}
\end{figure}

\begin{figure}[!t]
    \centering
    \includegraphics[width=\linewidth]{evaluation_protocol.png}
    \caption{Evaluation Protocol: Dataset branches into TRAIN (Model Training $\rightarrow$ Score Prediction $\rightarrow$ Evaluation) and TEST (Candidate Set 1 pos + 99 neg $\rightarrow$ Ranking).}
    \label{fig:evaluation_protocol}
\end{figure}

\end{document}
